\section{Rate-Complexity Bounds}\label{sec:bounds}
%==============================================================================

The threshold $C_0=1$ has an operational consequence: independent rate governs manual synchronization cost. This section states that consequence in the same finite, deterministic model used for the zero-error theorems.

\subsection{Cost Model}\label{sec:cost-model}

\begin{definition}[Modification Cost Model]\label{def:cost-model}
Let $\delta_F$ be a modification to fact $F$ in encoding system $C$. The \emph{effective modification complexity} $M_{\mathrm{effective}}(C,\delta_F)$ is the number of locations that must be edited manually to preserve correctness after applying $\delta_F$.
\end{definition}

Only manual edits are counted. Derived locations updated automatically by the host system contribute zero effective cost.

\subsection{Upper Bound at Unit Independent Rate}\label{sec:upper-bound}

\begin{theorem}[Rate-1 Upper Bound]\label{thm:upper-bound}
If $\text{DOF}(C,F)=1$, then
\[
M_{\mathrm{effective}}(C,\delta_F)=O(1).
\]
\end{theorem}
\leanmeta{\LH{BND1}}

\begin{proof}
When $\text{DOF}(C,F)=1$, there is exactly one independently writable source location for $F$. Updating that location is one manual edit; every other representation of $F$ is derived and is updated automatically. The number of manual edits therefore stays bounded by a constant independent of the number of visible encodings.
\end{proof}

\subsection{Lower Bound Above the Threshold}\label{sec:lower-bound}

\begin{theorem}[Above-Threshold Lower Bound]\label{thm:lower-bound}
If fact $F$ is encoded at $n$ independent locations, then
\[
M_{\mathrm{effective}}(C,\delta_F)=\Omega(n).
\]
\end{theorem}
\leanmeta{\LH{BND2}}

\begin{proof}
Each independent location can retain an outdated value unless it is edited directly. To restore exact consistency after changing $F$, each of the $n$ independent locations must therefore be synchronized manually. Hence the effective modification complexity grows at least linearly in $n$.
\end{proof}

\subsection{Finite Counting Converse}\label{sec:info-converse}

The update-cost lower bound is the operational face of the same information obstruction studied earlier: insufficient side information prevents exact recovery of the correct latent value.

\begin{theorem}[Finite Counting Converse]\label{thm:fano-converse}
Let $F$ range over $K$ possible latent values. Suppose exact recovery is attempted from a fixed observation transcript together with an $L$-bit side-information tag. Then
\[
K \le 2^L.
\]
Equivalently, exact zero-error recovery of $K$ ambiguous states requires at least $\log_2 K$ bits of side information.
\end{theorem}
\leanmeta{\LH{CIA3}}

\begin{proof}
An $L$-bit tag yields at most $2^L$ distinct labels. If the base observation transcript is fixed across the ambiguous states, exact recovery requires the tag assignment to distinguish all $K$ possibilities. Distinct states therefore need distinct tags, which is possible only when $K\le 2^L$.
\end{proof}

\begin{lemma}[Information-Constrained DOF Lower Bound]\label{lem:info-dof}
If the available side-information mechanism exposes at most $2^L$ distinguishable tags while the latent ambiguity class has size $K>2^L$, then no zero-error resolver can identify the true state from those tags alone. Any architecture that still guarantees exact correctness must therefore rely on additional independently accessible discriminating support, moving the system above the single-source regime.
\end{lemma}
\leanmeta{\LH{CIA4}}

\begin{proof}
Theorem~\ref{thm:fano-converse} rules out exact recovery from the available tag budget. If exact correctness is nevertheless required, the architecture must supplement the insufficient side-information channel with additional independently accessible information. In the present model, that means leaving the rate-$1$ regime and paying the synchronization cost of Theorem~\ref{thm:lower-bound}.
\end{proof}

\subsection{Worked Example}\label{sec:worked-mi}

Consider a fact with $K=4$ admissible values while the available side-information channel supplies only one bit. Since $2^1 < 4$, Theorem~\ref{thm:fano-converse} rules out zero-error recovery. The architecture must either enlarge the side-information alphabet or expose additional independent support. The theorem does not depend on software; the same logic applies to any host system whose observation transcript leaves four mutually compatible alternatives.

\subsection{Unified Converse Viewpoint}\label{sec:unified-converse}

The counting theorem is the basic form of the converse, but it is not the only one. In the deterministic finite model, the same budget obstruction can be restated in several classical languages.

\begin{theorem}[Unified finite converse, summary form]
Let $X$ be the latent source, let $Y$ denote the deterministic observation/side-information transcript, and let $\hat X$ be any deterministic decoder output. If $Y$ does not identify the true state exactly, then the same finite budget simultaneously:
\begin{enumerate}
\tightlist
\item lower-bounds zero-error ambiguity resolution and confusability-clique separation,
\item upper-bounds the conditional entropy $H(X\mid Y)$,
\item upper-bounds the decoder-output entropy $H(\hat X)$,
\item and appears as an output-entropy or KL gap relative to the ambient finite alphabets.
\end{enumerate}
These are different normal forms of one deterministic finite converse family rather than unrelated lemmas.
\end{theorem}

The zero-incoherence threshold identifies the unique exact-consistency regime; the converse family explains how departures from that regime reappear as finite information-budget obstructions.

\subsection{The Unbounded Gap}\label{sec:gap}

\begin{theorem}[Unbounded Gap]\label{thm:unbounded-gap}
The ratio between the modification cost of an architecture with $n$ independent encodings and the modification cost of a rate-$1$ architecture diverges:
\[
\lim_{n\to\infty}\frac{M_{\mathrm{DOF}>1}(n)}{M_{\mathrm{DOF}=1}}=\infty.
\]
\end{theorem}
\leanmeta{\LH{BND3}}

\begin{proof}
By Theorem~\ref{thm:upper-bound}, a rate-$1$ architecture has constant effective modification complexity. By Theorem~\ref{thm:lower-bound}, an architecture with $n$ independent encodings incurs linear cost in $n$. The ratio therefore grows without bound.
\end{proof}

The asymptotic statement has a simple practical reading: once a fact is duplicated across many independently writable locations, the cost of preserving exact consistency scales with the number of independent copies, not merely with the visibility of the fact.
